% Pursuing Stacks (cotes 134-2 to 134-8): archive cote 134-4, site page 7 (PDF page 8), typescript page 95.
% typescript. Draft compared: ps3.tex:316-382.
% AI-transcribed by Claude Opus 5.5 (claude-opus-5-5, adaptive thinking, effort high), Message Batches API, batch ps1 (msgbatch_01Rv5hCaNTtUt4JHx1yqoHPL), 2026-11-12; image review recorded in records/CHECK_DECISIONS_134_3_4.tsv.
\scanpage{134-4/7}
\mspage{(95)}

\uline{Question}  Let
\[ f:\ M^{\unsure{*}}\hand{\longrightarrow} M' \]
\margin{\diagram{\begin{tikzpicture}[x=1ex,y=1ex]\draw (0,0)--(.12,5)--(.1,10)--(.2,14);\end{tikzpicture}}}be a functor between modelizers (M,W) and (M',W'),  assume \xout{possibly} if
needed that f commute with  inverse limits,  or even has a left adjoint,
and that \xout{projective} inverse limits (and direct ones too,  as f\over[typed]{a}{o}r that !)
exist in M,M'.  Are there some natural conditions we\xout{x} can devise for
M and M'  (which should be satisfied for elementary modelizers and for
the basic modelizer (Cat)),  plus possibly some \over[typed]{m}{m}\over[typed]{i}{i}ld extra conditions on
f itself,  which will ensure that whenever f transforms weak equivalences
into weak equivalences,  or even only aspheric objects into aspheric ob-%
jects,  f is model-preserving i.e.  $W_M = f^{-1}(W_{M'}\ )$  \uline{and} the induced funct\unsure{or}
$H_f : H_M \hand{\longrightarrow} H_{M'}$  on the localizations is an equivalence of categories.

May\over[typed]{v}{b}e it's a silly question,  with pretty obvious negative answer
-  in any case,  I'll have to find out !  T\over{h}{h}e very first thing to check
is to see what happens in case of a functor
\[ i^{*}\ :\ M\hand{\longrightarrow} A^{\hand{\wedge}}\ , \]
where $A^{\hand{\wedge}}$ is a weak elementary modelizer,  and where M is either (Cat) or
another weak elementary modelizer $B^{\hand{\wedge}}$,  $i^{*}$ in any case being associated of
course to a functor
\[ i:\ A\hand{\longrightarrow} M\ \ , \]
with \inserted{a priori} no special requirement whatever on i.   In case  M = (Cat),  this
means lo\over[typed]{v}{o}king up in the end the question we have postponed for \inserted{quite} a while
now,  namely of how to rescue the ``key result'' of p.61,  when dropping the
assumption that the categories i(a)  (for a in A) have final objects. We
finally got a strong motivation for carrying through a generalization,  if
this is indeed feasible.

30.3.

47.     It had become clear that the most urgent thing to do now was to
come to a better understanding of test functors with values in (Cat),  when
dropping the assumption that the categories i(a) have final objects,  and
trying to replace this (if it should turn out that something \uline{is} needed in-%
deed)  by a kind of assumption which should make sense when (Cat) is re-%
placed by a more or less arbitrary modelizer M.  I spent a few hours pon-%
dering over the situation,  and it seems to me that\hand{\unsure{,}} in the case at least
when A is a \uline{strict} test category\hand{,} namely when $A^{\hand{\wedge}}$\hand{/} is totally aspheric,  the\unsure{-}%
re is now \unsure{a} \inserted[typed]{rather} complete understanding of the situation,  with a generalizati\unsure{-}%
on of the ``key result'' of p.61 which seems to be wholly satisfactory.

\margin{7}
